% ============================================================
\chapter{Strange Attractors}
\label{ch:attractors}
% ============================================================

\section{Introduction}

In 1971, David Ruelle and Floris Takens coined the term ``strange attractor'' to describe the geometric objects that govern the dynamics of chaotic systems. The Lorenz attractor, discovered in 1963, is the prototype: a butterfly-shaped set of fractal dimension $\approx 2.06$ on which trajectories wind without ever retracing their path. The H\'enon attractor, the R\"ossler attractor, and many others followed. What makes these objects ``strange'' is the coexistence of two contradictory behaviours: attraction (nearby orbits converge toward the attractor) and chaos (once on the attractor, orbits diverge exponentially). Geometrically, this coexistence produces a self-similar fractal structure, the result of repeated stretching and folding of phase space.

\begin{intuition}
A strange attractor works like a dough-kneading machine: it stretches nearby
orbits (creating sensitivity to initial conditions) and then folds them back
into a bounded volume (ensuring compactness). The result is an infinitely
fine layered structure --- a fractal.
\end{intuition}

\section{Formal Definitions}

\begin{definition}[Attractor]
\index{attractor}
Let $\varphi_t$ be a flow on $\R^n$. A compact invariant set $A$ is an
\textbf{attractor} if there exists an open neighbourhood $U \supset A$
(basin of attraction) such that:
\begin{enumerate}
    \item $\varphi_t(U) \subset U$ for all $t > 0$;
    \item $A = \bigcap_{t > 0} \varphi_t(U)$;
    \item $A$ is topologically transitive (contains a dense orbit).
\end{enumerate}
\end{definition}

\begin{definition}[Strange attractor]
\index{strange attractor!definition}
An attractor is \textbf{strange} if it has a fractal structure (non-integer
dimension) and exhibits sensitivity to initial conditions (at least one
positive Lyapunov exponent).
\end{definition}

\begin{attention}[title=\textbf{Strange $\neq$ chaotic}]
There exist strange non-chaotic attractors (fractal dimension but non-positive
Lyapunov exponents) and chaotic non-strange attractors (chaos on a smooth
torus). The two notions are logically independent.
\end{attention}

\section{The Lorenz Attractor}

\begin{definition}[Lorenz system]
\index{Lorenz system}
The \textbf{Lorenz system} (1963) is
\[
    \begin{cases}
    \dot{x} = \sigma(y - x), \\
    \dot{y} = \rho x - y - xz, \\
    \dot{z} = xy - \beta z,
    \end{cases}
\]
with classical parameters $\sigma = 10$, $\rho = 28$, $\beta = 8/3$.
\end{definition}

\begin{theorem}[Tucker, 2002]
\index{Tucker's theorem}
For the classical parameters $(\sigma, \rho, \beta) = (10, 28, 8/3)$, the
Lorenz system possesses a \textbf{robust strange attractor}. This result,
proved using rigorous computer-assisted methods, confirmed a conjecture open
since 1963.
\end{theorem}

\begin{proposition}[Properties of the Lorenz attractor]
The Lorenz attractor satisfies:
\begin{enumerate}
    \item \textbf{Dissipativity}: $\nabla \cdot f = -(\sigma + 1 + \beta) < 0$,
    so volumes contract exponentially.
    \item \textbf{Symmetry}: invariance under $(x, y, z) \mapsto (-x, -y, z)$.
    \item \textbf{Lyapunov spectrum}: $\lambda_1 \approx 0.91$,
    $\lambda_2 = 0$, $\lambda_3 \approx -14.57$.
    \item \textbf{Fractal dimension}: $D_{KY} \approx 2.06$.
\end{enumerate}
\end{proposition}

\begin{keyformulas}[title=\textbf{Equilibria of the Lorenz System}]
\begin{align*}
    &C_0 = (0, 0, 0), \quad \text{unstable for } \rho > 1. \\
    &C_{\pm} = (\pm\sqrt{\beta(\rho-1)}, \pm\sqrt{\beta(\rho-1)}, \rho-1),
    \quad \text{unstable for }
    \rho > \rho_H = \sigma\frac{\sigma + \beta + 3}{\sigma - \beta - 1}.
\end{align*}
\end{keyformulas}

\section{The R\"ossler System}

\begin{definition}[R\"ossler system]
\index{R\"ossler system}
The \textbf{R\"ossler system} (1976) is
\[
    \begin{cases}
    \dot{x} = -y - z, \\
    \dot{y} = x + ay, \\
    \dot{z} = b + z(x - c),
    \end{cases}
\]
with typical parameters $a = 0.2$, $b = 0.2$, $c = 5.7$.
\end{definition}

\begin{remark}
Unlike the two-lobed Lorenz attractor, the R\"ossler attractor has a single
band structure with a simpler folding mechanism. It is often used as the
minimal model for single-band chaos.
\end{remark}

\begin{proposition}
For standard parameters, the R\"ossler system exhibits a period-doubling
cascade as $c$ increases: periodic orbit ($c \approx 2.5$), doubling
($c \approx 3.5$), chaos ($c \approx 4.2$), periodic window ($c \approx 5.2$),
fully developed chaos ($c \approx 5.7$).
\end{proposition}

\section{The H\'enon Map}

\begin{definition}[H\'enon map]
\index{H\'enon map}
The \textbf{H\'enon map} is the planar diffeomorphism
\[
    H(x, y) = (1 - ax^2 + y,\; bx),
\]
with classical parameters $a = 1.4$, $b = 0.3$.
\end{definition}

\begin{proposition}[Properties of the H\'enon map]
\begin{enumerate}
    \item The Jacobian is constant: $\det DH = -b$. The map contracts areas
    by a factor $\abs{b} = 0.3$.
    \item The attractor has fractal dimension $D \approx 1.26$.
    \item Lyapunov exponents are $\lambda_1 \approx 0.42$ and $\lambda_2 \approx -1.62$.
    \item The attractor has a Cantor structure in the transverse direction.
\end{enumerate}
\end{proposition}

\begin{theorem}[Benedicks-Carleson, 1991]
\index{Benedicks-Carleson theorem}
For a set of parameters $(a, b)$ of positive Lebesgue measure with $b$ small,
the H\'enon map possesses a strange attractor carrying a unique ergodic SRB
(Sinai-Ruelle-Bowen) measure.
\end{theorem}

\section{Fractal Dimension}

\begin{definition}[Box-counting dimension]
\index{dimension!box-counting}
Let $A \subset \R^n$ be bounded. The \textbf{box-counting dimension} is
\[
    D_0 = \lim_{\epsilon \to 0} \frac{\ln N(\epsilon)}{-\ln\epsilon},
\]
where $N(\epsilon)$ is the minimum number of boxes of side $\epsilon$ needed
to cover $A$.
\end{definition}

\begin{definition}[Hausdorff dimension]
\index{dimension!Hausdorff}
The \textbf{Hausdorff dimension} of $A$ is
\[
    D_H = \inf\left\{ d \geq 0 : \mathcal{H}^d(A) = 0 \right\},
\]
where $\mathcal{H}^d(A) = \lim_{\epsilon \to 0} \inf\left\{
\sum_i (\mathrm{diam}\, U_i)^d : A \subset \bigcup_i U_i,\;
\mathrm{diam}\, U_i < \epsilon \right\}$.
\end{definition}

\begin{proposition}[Dimension inequality]
For any bounded set $A \subset \R^n$:
\[
    D_H(A) \leq \underline{D}_0(A) \leq \overline{D}_0(A).
\]
\end{proposition}

\begin{example}[Triadic Cantor set]
The triadic Cantor set has dimension
$D_H = D_0 = \frac{\ln 2}{\ln 3} \approx 0.631$.
At step $n$, there are $N(\epsilon) = 2^n$ intervals of length
$\epsilon = 3^{-n}$, giving $D_0 = \lim \frac{n\ln 2}{n\ln 3}$.
\end{example}

\begin{definition}[Correlation dimension]
\index{dimension!correlation}
The \textbf{correlation dimension} is
\[
    D_2 = \lim_{\epsilon \to 0} \frac{\ln C(\epsilon)}{\ln\epsilon},
\]
where $C(\epsilon) = \lim_{N \to \infty} \frac{1}{N^2}\#\{(i,j) :
\norm{x_i - x_j} < \epsilon\}$ is the Grassberger-Procaccia correlation
integral.
\end{definition}

\section{Strange Attractors in Physical Systems}

\begin{example}[Rayleigh-B\'enard convection]
The Lorenz system is a truncation of the Navier-Stokes equations for
Rayleigh-B\'enard convection between two heated horizontal plates. The
parameter $\rho$ corresponds to the normalised Rayleigh number.
\end{example}

\begin{example}[Chua's circuit]
\index{Chua's circuit}
\textbf{Chua's circuit} is a simple electronic circuit (capacitors, inductor,
nonlinear resistor) that produces a double-scroll strange attractor. It is the
first physical system for which chaos was rigorously proved.
\end{example}

\begin{codeblock}[title=\textbf{Visualising the Lorenz attractor}]
\begin{minted}{python}
import numpy as np
from scipy.integrate import solve_ivp
import matplotlib.pyplot as plt

def lorenz(t, state, sigma=10, rho=28, beta=8/3):
    x, y, z = state
    return [sigma*(y - x), rho*x - y - x*z, x*y - beta*z]

sol = solve_ivp(lorenz, [0, 100], [1, 1, 1],
                max_step=0.01, t_eval=np.linspace(0, 100, 100000))
fig = plt.figure(figsize=(10, 8))
ax = fig.add_subplot(111, projection='3d')
ax.plot(sol.y[0], sol.y[1], sol.y[2], lw=0.3, color='darkblue')
ax.set_xlabel('$x$'); ax.set_ylabel('$y$'); ax.set_zlabel('$z$')
ax.set_title('Lorenz Attractor')
plt.savefig("lorenz_attractor.pdf")
\end{minted}
\end{codeblock}

\section{Exercises}

\begin{exercise}[Dissipativity]
Show that the Lorenz system is dissipative and that any initial volume $V_0$
evolves as $V(t) = V_0 e^{-(\sigma + 1 + \beta)t}$. Deduce the existence
of a zero-Lebesgue-measure attractor.
\end{exercise}

\begin{exercise}[Box-counting dimension]
Compute the box-counting dimension of $\{0\} \cup \{1/n : n \in \N^*\}$
and of the Menger sponge.
\end{exercise}

\begin{exercise}[H\'enon attractor]
Write a program that plots the H\'enon attractor for $a = 1.4$, $b = 0.3$
by iterating $10^6$ times. Estimate its correlation dimension using the
Grassberger-Procaccia method.
\end{exercise}

\begin{exercise}[Poincar\'e section]
For the R\"ossler system, plot the Poincar\'e section in the plane
$z = z_{\max}/2$ and show that it resembles a one-dimensional map.
\end{exercise}

\begin{exercise}[Structural stability]
Is the Lorenz attractor structurally stable? Discuss by comparing with
Smale's horseshoe and justify the notion of a robust attractor (in the sense
of Morales-Pac\'ifico-Pujals).
\end{exercise}
